\thispagestyle{empty}
\begin{flushleft}
{\footnotesize УДК~004.056:004.492.3(075.8)\par}
\end{flushleft}

\vspace{2mm}

{\footnotesize\noindent
\textbf{Соколов~В.} Технології протидії злоякісному програмному коду: навчальний посібник. --- Київ, 2026. --- 17~розділів.\par}

\vspace{4mm}

{\footnotesize\noindent
Навчальний посібник присвячено технологіям виявлення, аналізу та протидії злоякісному програмному коду. Систематично розглянуто поняттєвий апарат і таксономію шкідливого програмного забезпечення, нормативно-правове й стандартизаційне забезпечення, теоретичні та математичні основи виявлення, моделювання загроз і ланцюги ураження, внутрішню будову та життєвий цикл шкідливого коду. Докладно висвітлено статичний і динамічний аналіз, сигнатурне та поведінкове виявлення, техніки ухилення й обфускації, застосування машинного навчання, захист кінцевих точок і телеметрію, інженерію виявлення, розвідку загроз і атрибуцію, протидію цілеспрямованим сталим загрозам (APT) та полювання за загрозами. Окремо розглянуто захист спеціалізованих середовищ (промислові системи, Інтернет речей, мобільні платформи, хмару), атаки на ланцюг постачання, реагування на інциденти й цифрову форензику, а також організаційне управління та стратегію кіберстійкості. Матеріал спирається на актуальні міжнародні стандарти й наукові джерела та враховує досвід протидії кіберзагрозам в Україні.\par}

\vspace{4mm}

{\footnotesize\noindent
Для здобувачів третього (освітньо-наукового) рівня вищої освіти спеціальності 125/F5 «Кібербезпека та захист інформації», а також науковців, фахівців із кібербезпеки та всіх, хто цікавиться технологіями протидії шкідливому програмному забезпеченню.\par}

\vfill

{\footnotesize\noindent
\emph{Матеріали посібника мають виключно освітній та дослідницький характер і спрямовані на захист інформаційних систем. Наведені відомості не містять працездатних зразків шкідливого коду й призначені для підготовки фахівців із кіберзахисту.}\par}

\vspace{3mm}

{\footnotesize\noindent \textcopyright{} Соколов~В., 2026\par}
